\documentclass[11pt]{article}
\usepackage[margin=2.2cm]{geometry}
\usepackage{amsmath,amssymb,listings}
\usepackage{booktabs}
\usepackage{xurl}
\usepackage{xcolor}
\usepackage[colorlinks=true,linkcolor=blue!50!black,citecolor=blue!50!black,urlcolor=blue!50!black]{hyperref}

\usepackage{orcidlink}

\lstset{
    language=Python,
    basicstyle=\ttfamily\footnotesize,
    keywordstyle=\color{blue!60!black}\bfseries,
    stringstyle=\color{green!45!black},
    commentstyle=\color{gray}\itshape,
    breaklines=true,
    frame=single,
    framerule=0.3pt,
    xleftmargin=1em,
}
\input{gen/leanos_facts.tex}

\title{LeanOS: a kernel written to be proved by LEAN4\\
\large from a proof-checked loader to a proof-checked memory map}
\author{Brent S. Hartshorn \orcidlink{0009-0004-2853-655X}
        \small (\url{brenthartshorn@proton.me}) 
\\
\small
\url{https://github.com/brentharts/RosettaMath}
\quad
\url{https://github.com/brentharts/crust}
}
\date{September 14, 2026}

\begin{document}
\maketitle

\begin{abstract}
LeanOS, an operating system micro-kernel written to be proved by LEAN4,
built so far: an ELF validator that decides before it dereferences, with a loader
that now refuses what it used to run; the region list LeanOS is founded on,
whose founding sentence~--- no address is reachable by two owners~--- is now
a theorem; threads as records over that list; a per-thread bump allocator;
a loader that makes the ELF theorem and the memory theorem one check; and
the switcher's decisions in C, lifted from the compiled binary and proved
from the lift.  Every LEAN theorem is checked by \texttt{lean4.py} and again by
Lean~4 with no axioms.  We record what the theorems
do not say, and the additions to the proof kernel that were needed to say
this much.
\end{abstract}

\section{Introduction}

Three things in this tree are checked by a proof kernel.  A SIMD contract is
certified at compile time by \texttt{lean4.py}, a Calculus of Constructions
kernel written in Python and named after Lean~4 rather than being it.  A
Python \emph{model} of the kernel is proved by \texttt{hoare.py} and
exported to Lean~\leanVersion{}, which checks it independently.  And
\texttt{ilproof.py} lifts a compiled function's IL back into the fragment
\texttt{hoare.py} reads, so that a theorem can be about code the compiler
actually saw rather than a transcription of it.

The lift is where the two kernels meet, and its tally is the reason for
this document.  Of the \liftTotal{} functions in \texttt{crustos/kernel.c},
the lift reaches \liftReached{}.  Of the \liftRefused{} it refuses,
\liftMemory{} are refused for one reason: \texttt{ReadAt}, \texttt{AddrOf}
or \texttt{SetAt}~--- a memory access through a pointer the lift has no
way to name.  The refusals are not spread across the kernel's features.
They are one wall.

The reading of that wall which LeanOS is built on: \emph{a pointer the
lift cannot name is a region the kernel never wrote down.}  If every region
the kernel touches were an entry in a list laid out at build time~--- base,
size, owner~--- then a load from one would be a field read on a record the
lift can already express, not a \texttt{ReadAt} it must refuse.  The
\liftMemory{} refusals are the count of pointers such a list would have
replaced.

\section{The loader, first}

The loader was chosen as the first thing to write in the new style because
its correct behaviour is a property of \emph{headers}, and headers are a
list of records.  Deciding things about lists is what rpython is for in
this tree, and it is what the proof side can already reason about.

\subsection{The validator}

\texttt{crustos/elfcheck.py} is four rpython functions.  It receives the
\texttt{PT\_LOAD}s as two parallel lists in file order, with the entry
point and the register class from the \texttt{CRUSTOS} note, and it answers
one or zero.  It never reads memory it did not receive as a list.

\lstinputlisting[firstline=34,lastline=62]{\crustDir/crustos/elfcheck.py}

Three conventions make the model and the code the same function.  A
result is 1 or 0, never a negative sentinel, because the model is over
\texttt{Nat}.  A loop counts up, so its variant is \texttt{len - i}.  And
\texttt{reg\_class\_ok}'s guard is written \texttt{<=} rather than
\texttt{>}, so that the guard \emph{is} the theorem: an accepted class is
\texttt{<= 3} because that is what was tested, and the kernel settles it by
splitting the guard rather than by a totality lemma about \texttt{leb} that
the prelude does not have.

\subsection{What is proved, and by which kernel}

RosettaMath's \texttt{elfcheck\_eq.py} is the same four functions in the
dialect \texttt{read\_procedure} compiles.  Its \elfTheorems{} theorems are
proved by \texttt{lean4.py} in \elfBuildSeconds{}s, exported as
\elfLeanLines{} lines of Lean, and \elfLeanVerdict{} by Lean~4 in
\elfLeanSeconds{}s with \elfLeanNoAxioms{} of \elfTheorems{} reported free
of axioms:
\begin{align*}
\texttt{reg\_class\_sized} &: \forall\,\mathit{cls},\ \texttt{reg\_class\_ok}\ \mathit{cls} = 1 \to \mathit{cls} \le 3\\
\texttt{accept\_sized} &: \forall\,v\,m\,e\,\mathit{cls},\ \texttt{accept\_image}\ v\,m\,e\,\mathit{cls} = 1 \to \mathit{cls} \le 3\\
\texttt{loads\_ordered\_bounded} &: \forall\,v\,m,\ \texttt{loads\_ordered}\ v\,m \le 1\\
\texttt{entry\_in\_load\_bounded} &: \forall\,v\,m\,e,\ \texttt{entry\_in\_load}\ v\,m\,e \le 1
\end{align*}
The second composes with \texttt{elf\_regs\_bounded}, the one theorem
already lifted from \texttt{kernel.c} through the real front end and
accepted by Lean: an image the loader accepts has a class the scheduler
sizes, and that sizing is \texttt{<= 23}.  It is the first theorem in the
tree that spans two functions of the shipped kernel.

\subsection{The loader before and after}

\texttt{elf\_load\_path} now builds the two lists after its first pass
over the program headers, with the same \texttt{memsz != 0} filter the
mapping loop uses, and returns \texttt{-6} if \texttt{accept\_image} says
no, before anything is allocated.  Table~\ref{tab:loader} is a
ShivyCX-built ELF and four copies of it with patched headers, through the
loader as it was and as it is.

\begin{table}[h]\centering
\begin{tabular}{lcc}
\toprule
& before & after\\
\midrule
well-formed & \texttt{rc=0} & \texttt{rc=0}\\
second load begins inside the first & \texttt{rc=0} & \texttt{rc=-6}\\
loads in descending order & \texttt{rc=0} & \texttt{rc=-6}\\
entry in no load, 1\,MB past the image & \texttt{rc=0} & \texttt{rc=-6}\\
entry exactly at the end of the last load & \texttt{rc=0} & \texttt{rc=-6}\\
entry at \texttt{end - 1} & \texttt{rc=0} & \texttt{rc=0}\\
\bottomrule
\end{tabular}
\caption{\texttt{elf\_load\_path} on real files.  Before the validator,
\texttt{elf\_run\_guest\_fn} would have jumped to the entry in row four.}
\label{tab:loader}
\end{table}

\subsection{How the validator is checked}

Five ways, and the first four are the discipline
\texttt{test\_crustos\_model.py} already applies to the scheme layer.
\emph{Behaviour}: the rpython run as Python and the compiled model terms
evaluated through the kernel agree over a corpus.  \emph{Shape}: the two
sources have the same control-flow skeleton, so a corpus cannot miss a
branch.  \emph{Coverage}: every exported function is modelled.
\emph{Lean}: the export is accepted with no axioms.  And two the scheme
layer does not have.  \emph{Compiled}: the same file through ShivyCX by
\texttt{\#include}, as \texttt{kernel.c} takes it, against itself run as
Python.  \emph{End to end}: the loader on the files of
Table~\ref{tab:loader}.

The tamper that swaps \texttt{<} for \texttt{<=} on the entry bound~---
the loader's off-by-one, which lets an entry sit one byte past its
segment~--- is caught by the behaviour corpus on three rows.

\section{The founding rule}\label{sec:founding}

\textbf{Memory is data before it is memory.}  Every region LeanOS will
touch is an entry in one list: base, size, owner.  Disjointness of the
list is the theorem everything else rests on, because a thread allocated
from its own region has no address in its registers or its slots that
lands in another's.  The register partition of
\texttt{BAREMETAL\_THREADS.md}~--- two threads, one core, disjoint
footprints, a switcher that saves what the partition says~--- is not a
special case of this but its first instance.  LeanOS makes the same move
for memory, for the same reason: so that the partition is a fact the
compiler computed and the kernel can check.

\subsection{The region list}

\texttt{leanos/memmap.py} is written the way \texttt{elfcheck.py} is, and
deliberately: \texttt{regions\_disjoint} is \texttt{loads\_ordered} with an
owner column, so the loop theorem each owes is one theorem.

\lstinputlisting[firstline=22,lastline=41]{\crustDir/leanos/memmap.py}

\texttt{region\_of} is the access rule.  A thread may touch an address only
if it says 1, and with disjointness in hand, 1 for one owner is 0 for every
other.  Compiled by ShivyCX, laid out as the two-thread partition~--- kernel
tables, two stacks, two slot regions~--- it finds an address in thread~1's
stack, denies it to thread~2 and to the kernel, and reports the layout
non-disjoint the moment thread~1's stack is grown one byte into its slots.

\subsection{What is proved}

\texttt{memmap\_eq.py} proves \memTheorems{} theorems in
\memBuildSeconds{}s; Lean~4 \memLeanVerdict{} the \memLeanLines{}-line
export in \memLeanSeconds{}s, \memLeanNoAxioms{} of \memTheorems{} free of
axioms.  Five say a function decides.  The rest are a chain, and the chain
ends at the founding sentence of \S\ref{sec:founding} stated verbatim:
\begin{align*}
\texttt{regions\_disjoint\_ordered} &: \texttt{regions\_disjoint}\ b\,s = 1 \to \texttt{ordered\_prefix}\ b\,s\ (\mathrm{len}\ b)\\
\texttt{regions\_pairwise\_disjoint} &: \texttt{ordered\_prefix}\ b\,s\,n \to j<k<n \to b_j+s_j \le b_k\\
\texttt{contains\_unique} &: \text{two positions holding the same address are one position}\\
\texttt{region\_index\_found} &: \text{the index returned is owned by the asker and holds the address}\\
\texttt{region\_of\_unique} &: \texttt{region\_of}\ w_1\,a = 1 \to \texttt{region\_of}\ w_2\,a = 1 \to w_1 = w_2
\end{align*}
The first of these,
\[
\texttt{regions\_disjoint}\ b\,s = 1 \;\to\; \texttt{ordered\_prefix}\ b\,s\ (\mathrm{len}\ b),
\]
says that if the checker answers 1, every consecutive pair of regions is
in order, and \texttt{regions\_checked\_is\_len} is the equation the exit
condition yields.  It did not need a ``for every element'' primitive.  The
specification \texttt{ordered\_prefix} is written with the lowering's own
guard expression~--- the same \texttt{ltb}, \texttt{nth}, \texttt{add} and
\texttt{sub i 1}~--- so that once a guard is split the specification
\emph{computes} rather than needing a lemma to connect it to the code.  And
the code carries no \texttt{end} variable, recomputing the previous end from
\texttt{i}, because carrying it would have made the invariant an equation
between a variable and what it was meant to hold, and an equation is the one
thing the fragment cannot rewrite with.  The guard is the theorem, and this
time it reached into the shape of the loop.

The last link needed a change to the shipped code, and it is a design
rule: \texttt{region\_of} returned a Bool, and a proof about \emph{which}
region was found cannot be written against a Bool.  It is now
\texttt{region\_index}, which returns the witness or \texttt{len(bases)},
with \texttt{region\_of} a two-line guard over it~--- exactly as
\texttt{scheme\_of} returns an index and not a flag.  Write the code so the
proof can see what it did.

\section{On the region list: threads, allocation, admission}

Milestone~0 was expensive because it built the machinery; milestones 1 to 3
were cheap because they only used it.  Each is a file in \texttt{leanos/},
compiled by ShivyCX together with \texttt{memmap.py} in one C unit, checked
the five ways of \S2.4, and each theorem is a corollary of
\texttt{region\_of\_unique} or of a guard.

\paragraph{Threads as records.}  \texttt{threads.py}: a thread is an owner
(\texttt{tid + 1}; the kernel is 0) and a stack pointer.  \texttt{sp\_ok} is
the access rule applied to the stack, and \texttt{sp\_after\_push} moves the
pointer down by $n$ \emph{if the result is still in the thread's region},
otherwise leaves it~--- a stack overflow is a decision the kernel makes with
the region list in hand, not a fault it discovers afterwards.  \thrTheorems{}
theorems, \thrLeanVerdict{} by Lean in \thrLeanSeconds{}s:
\begin{align*}
\texttt{push\_keeps\_sp\_ok} &: \texttt{sp\_ok}\ t\,\mathit{sp} = 1 \to \texttt{sp\_ok}\ t\ (\texttt{sp\_after\_push}\ \ldots\ \mathit{sp}\ n) = 1\\
\texttt{sp\_no\_cross} &: \texttt{sp\_ok}\ a\,\mathit{sp} = 1 \to \texttt{sp\_ok}\ b\,\mathit{sp} = 1 \to a = b
\end{align*}
The first is ``the guard is the invariant, so preservation is the guard'':
the leaf that returns $\mathit{sp}-n$ is settled by the very guard that
admitted it.  Neither needed a new lemma.

\paragraph{The per-thread allocator.}  \texttt{alloc.py}: no shared heap.
\texttt{bump} moves \texttt{used} up a region the thread owns if the request
fits, else leaves it.  \alcTheorems{} theorems, \alcLeanVerdict{} in
\alcLeanSeconds{}s: \texttt{bump\_bounded} ($\mathit{used}\le\mathit{size}$
is kept), \texttt{bump\_monotone} (allocation never gives back),
\texttt{slot\_in\_bounds} (the slot handed out is inside the region).  This
is region-based allocation in Gay's sense~\cite{gay2001}, with the region
structure fixed at build time rather than inferred, and the theorems about
the structure rather than about a runtime.

\paragraph{Admission.}  \texttt{loader.py}: \texttt{admit} is
\texttt{accept\_image} on the headers and then \texttt{regions\_disjoint} on
the region list with the loads appended under the guest's owner, both guards
written \texttt{== 1} so each guard is the theorem about it.
\ldrTheorems{} theorems, \ldrLeanVerdict{} in \ldrLeanSeconds{}s;
the one that is the milestone is the composition
\[
\texttt{admit\_ordered} : \texttt{admit} = 1 \to \texttt{ordered\_prefix}\ (\texttt{extended}\ b\,v)\ (\texttt{extended}\ s\,m)\ (\mathrm{len}\ \ldots),
\]
which is \texttt{regions\_disjoint\_ordered} on the grown list: every
theorem of \S3 holds with the guest in it.  The ELF theorem and the memory
theorem are one because the loader makes them one check.

\section{The switcher, from the binary}

\texttt{leanos/switch.c} is the first C in LeanOS and the first thing
proved from the compiled binary rather than from source.  It is the
switcher's four decisions~--- which side owns a register, where it is
saved, how many a side saves, who runs next~--- as integer functions of
the partition of \texttt{BAREMETAL\_THREADS.md} (left $=$ x19--x24, right
$=$ x25--x28).  \texttt{ilproof.py} lifts all \swFunctions{} from the IL the
compiler produced; the binary agrees with the lift over \swRegs{} register
numbers; and \swTheorems{} theorems proved from the lifted source are
\swLeanVerdict{} by Lean in \swLeanSeconds{}s, \swLeanNoAxioms{} free of
axioms.  The two that matter say the footprints are disjoint,
\[
\texttt{reg\_side}\ r = 0 \to r < 25 \qquad\qquad \texttt{reg\_side}\ r = 1 \to \neg(r<25),
\]
and \texttt{saves}\ 0 $= 6$ is the six-register footprint the measured
case reports.  \texttt{save\_slot}\ r for $r<19$ is negative in C and 0
over Nat; the test states that as the one column where model and machine
part.

What is not here is the saving itself: the stores and loads of register
values into the thread's slots are memory access, which the lift refuses.
That is the wall of \S1, and the region list was built so that each slot
is a field of a region the list names.  Lifting \texttt{SetAt} and
\texttt{ReadAt} against the list is the next compiler step.  A second
compiler item: the contract front end takes \texttt{len(x)} bounds only,
so \texttt{assert split >= 19} on a parameterised partition is not yet a
contract, and the partition is constants.

\section{What it cost the prover}
\label{sec:prover}

Additions to \texttt{hoare.py}, each general rather than fitted to the
case that exposed it.  The prelude went from 75 to \preludeDecls{}
declarations, every one a hand-written kernel term: totality, antisymmetry
and irreflexivity of \texttt{leb}; \texttt{eqb} as an equation and back;
\texttt{sub\_zero}, \texttt{le\_add\_right}, \texttt{add\_le\_add\_left};
and the Bool lemmas that read a decided disjunction.  Each is an induction
whose every case is computation, the hypothesis, or \texttt{absurd}.

\paragraph{A guard written in terms of another.}
\texttt{by\_every\_bool} settles a chain of guards by splitting each and
computing.  It split the \emph{outermost} first.  In \texttt{accept\_image}
the outermost guard is \texttt{eqb (reg\_class\_ok cls) 0}, and
\texttt{reg\_class\_ok} is itself an \texttt{ite} on \texttt{leb cls 3};
splitting the outer one as an atom left \texttt{leb cls 3} open everywhere
else it occurred.  It now splits the innermost open guard first.  And a
guard that \emph{computes} to a literal after earlier splits is written in,
not split: splitting it asks for a proof of the branch computation rules
out, and there is none.

\paragraph{A postcondition about an early return.}
\texttt{desugar\_returns} lowers a \texttt{return} inside a loop into a
first-wins accumulator \texttt{\_return\_value}, and reserved that name
against the body~--- including against an
\texttt{assert invariant(...)} that merely \emph{read} it.  So no
postcondition about a value returned early from a loop was provable: not
\texttt{loads\_ordered}'s, not \texttt{regions\_disjoint}'s, not
\texttt{scheme\_of}'s stated \texttt{result <= len(names)}.  The name is
now reserved against assignment.  With that, the recipe is:
entry by splitting the pre-loop guards, since the accumulator's seed is an
\texttt{ite} over them; preservation by splitting the body's guards and
then \texttt{\_returned}, where a branch that assigned a literal computes
and a branch that did not is the hypothesis; the variant by
\texttt{sub\_lt}; exit by splitting the final \texttt{\_returned}.
\texttt{early\_return\_bound} is that recipe as one call, and
\texttt{region\_of}~--- three nested guards, a different state
shape~--- went through it unchanged.

\paragraph{One weak-head step.}
After a pass, the claim arrives as projections of the state tuple:
\texttt{fst}/\texttt{snd} over \texttt{mk}.  \texttt{normalize} opens
\texttt{leb} and \texttt{ite} into recursors there is nothing left to match
against; unfolding \texttt{fst} and \texttt{snd} makes the term larger,
since each is \texttt{Prod.rec} under a lambda.  \texttt{reduce\_projections}
is the twelve lines that were missing.

\paragraph{Two equations, and two inductions.}
\texttt{holds\_orb\_left}, $\mathrm{Holds}\ a \to \mathrm{Holds}\ (\mathrm{orb}\ a\ b)$, is
the first place a loop proof needed an equation rather than a case
split~--- the hypothesis is about a variable, the goal about a term built
from it~--- and is ten lines by \texttt{Eq.ind} once \texttt{eq\_symm}
exists.  \texttt{ltb\_false\_leb} and \texttt{leb\_antisymm} are inductions
on the first argument with a case split on the second, each case
computation, the hypothesis, or \texttt{absurd}; together they turn ``the
counter is $\le$ the bound and not $<$ it'' into the counter being the
bound.

\paragraph{A split that keeps its evidence.}
\texttt{by\_bool} abstracts a guard out of the goal and applies
\texttt{Bool.ind}; a branch learns the goal has \texttt{true} written in,
not that the guard \emph{is} true.  That suffices when the goal's own
occurrences of the guard are what is decided.  It fails the moment a
branch must prove something about an occurrence that appears only after the
\texttt{ite} reduces: a search returning \texttt{i} must show the guards
hold \emph{at i}, and those copies were never substituted.
\texttt{by\_bool\_with\_evidence} carries the equation in its motive, and
\texttt{bound\_by\_ites} retries with it when the plain split will not
close.  Getting it right took four turns and found four bugs, one of them a
soundness hole of my own: a hypothesis about a variable offered as a proof
of a claim about \texttt{true}, because the comparison had been made after
substitution and the term had not been transported.  What settled it was
not another hypothesis but a tool~--- type-check the built term in the
leaf's own scope and bisect to the smallest failing piece.  This is the
kind of invariant Koenig's inference targets~\cite{koenig2022}: quantified
over positions, with the separation between regions doing the work.  Here
every invariant is written by hand, and the quantified ones were the
expensive ones.

\paragraph{And one that was not needed: division.}
\texttt{addr >> A::PAGE\_SHIFT} lifts as \texttt{addr // 4096}, which on an
unsigned value is the same number~--- no truncation, no wrap.  \texttt{divb}
in the prelude is \texttt{modb}'s counter incrementing where \texttt{modb}
resets, and Lean accepts it.  The lift reaches the kernel's two
\texttt{frame\_of} functions.  What it cannot yet do is prove
\texttt{frame\_of(8192) == 2}, which should be settled by computation and
is not: \texttt{normalize} opens the definition under its binder, where
\texttt{divb addr 4096} has a free variable and cannot accelerate, and
\texttt{eqb \_ 4096} unfolds \texttt{leb} four thousand levels deep.  Closed
evaluation is fine~--- \texttt{divb 131072 65536} is 0.66s~--- the
unfolding order is not, and that is a \texttt{lean4.py} matter.

\section{What is claimed, and what is not}

A theorem about a model is a theorem about a model.  The tests that run
the compiled code over the same corpus are what make it a theorem about
the kernel, and they are corpus tests: a corpus can miss an input.  The
shape test says a model cannot miss a branch, which is narrower and harder
to fake.

The validator and the region list are proved from \emph{source}.  Of the
validator's four functions, the lift reaches one from the compiled IL,
\texttt{reg\_class\_ok}; the three with loops are refused, because a list
index is a \texttt{ReadAt} at the IL level.  So the theorems here are about
what \texttt{read\_procedure} read, and the compiled-agreement test is the
bridge to what ShivyCX emitted.  Lifting list access is the same wall as
lifting the kernel's pointers, and the region list is the plan for both.

Addresses are \texttt{i64} in the code and \texttt{Nat} in the model;
\texttt{vaddr + memsz} wraps past $2^{63}$ in one and not the other.  The
model's numerals are unary, so a page-scale address is a term thousands
deep; the corpora use small numbers, and the compiled leg runs page-scale
and \texttt{0x7f...} cases because C does not care.

Nothing here is a proof about machine code.  The distance from IL to
what ShivyCX emits is translation validation, and it is the far end of
this road.

\section{Results: The Milestones (in proof order)}

\begin{enumerate}
\item \textbf{The region list, and its theorem.}  Done, to the founding
  sentence verbatim.
\item \textbf{Threads as records.}  Done.  \texttt{all\_sps\_ok}'s loop
  theorem is not stated.
\item \textbf{The per-thread allocator.}  Done.  The owner-level version of
  \texttt{slot\_in\_bounds} waits on the completeness of
  \texttt{region\_index}.
\item \textbf{The loader.}  Done.  ``The guest can reach its own entry'' is
  tested, and waits on the same completeness lemma.
\item \textbf{The switch.}  The decisions, done, from the binary.  The
  saving waits on memory access lifted against the region list.
\item \textbf{Syscalls.}  Each a function from a thread record to a
  thread record, which is what \texttt{tick} and \texttt{sched} already
  are~--- now about a \texttt{Context} that ships.  Not started.
\end{enumerate}

Two items block the last steps of three milestones: the completeness of
\texttt{region\_index}~--- if some region matches, it returns one~--- which
is an induction over positions of the kind \S\ref{sec:prover} now has the
tools for; and memory access in the lift, which is compiler work.

This draft continues the line of~\cite{hartshorn2026a}, which set out the
pipeline from \LaTeX{} theorems to compiler decisions, and
of~\cite{hartshorn2026b}, which had a second kernel agree with the first;
here the second kernel agrees about a kernel of the other kind.

\appendix
\section{The Generated LEAN4 Proofs}

\begin{enumerate}
\item alc.lean \url{https://github.com/brentharts/RosettaMath/wiki/alc.lean}
\item elf.lean \url{https://github.com/brentharts/RosettaMath/wiki/elf.lean}
\item ldr.lean \url{https://github.com/brentharts/RosettaMath/wiki/ldr.lean}
\item mem.lean \url{https://github.com/brentharts/RosettaMath/wiki/mem.lean}
\item sw.lean \url{https://github.com/brentharts/RosettaMath/wiki/sw.lean}
\item thr.lean \url{https://github.com/brentharts/RosettaMath/wiki/thr.lean}
\end{enumerate}


\begin{thebibliography}{9}
\bibitem{gay2001} David Gay.  \emph{Memory Management with Explicit
  Regions}.  PhD thesis, University of California, Berkeley, 2001.
  \url{https://theory.stanford.edu/~aiken/publications/theses/gay.pdf}
\bibitem{koenig2022} Jason Koenig.  \emph{Invariant Inference via
  Quantified Separation}.  PhD thesis, Stanford University, 2022.
  \url{https://theory.stanford.edu/~aiken/publications/theses/koenig.pdf}

\bibitem{hartshorn2026a} Brent Hartshorn.  ``LEAN Proofs Small Enough to
  Read: Kernel Contracts from \LaTeX{} Theorems to Compiler Decisions.''
  ai.viXra.org:2609.0019, 2026.
  \url{https://ai.vixra.org/pdf/2609.0019v1.pdf}
\bibitem{hartshorn2026b} Brent Hartshorn.  ``A Second Kernel Agrees:
  `LEAN Proofs Small Enough to Read', Checked by Lean4 and Read Four
  Ways.''  SSRN, 2026.  \url{https://dx.doi.org/10.2139/ssrn.7443439}
\end{thebibliography}

\end{document}
